\section{Bounded Change and the Safety Decision}
\label{sec:bounded_safety}

\subsection{Stage five, applying the change}
\label{sec:composition}
The admitted gain is the raw proposal modulated by the allocation and the strength map, then bounded,
\begin{equation}
\boldsymbol{\alpha} = \mathrm{clamp}\big(\boldsymbol{\Delta} \odot \mathbf{a} \odot \boldsymbol{\lambda},\, -0.20,\, 0.20\big).
\label{eq:alpha}
\end{equation}
A gate suppresses any change in the dark background,
\begin{equation}
\mathbf{g} = \sigma\big(20\,(\mathbf{b} - q_{0.15})\big),
\label{eq:gate}
\end{equation}
where $q_{0.15}$ is the fifteenth percentile of the backbone output of that image. Using a percentile
rather than a fixed level makes the gate work across scanners with different brightness. A second form
of the gate, used where the selection of Section~\ref{sec:weight_selection} chose it, is a soft step at
the per image Otsu threshold followed by a maximum filter of fixed radius, so that the gate is
identically zero beyond that radius from tissue.

The gated gain is applied multiplicatively to give the change $\mathbf{v}$, with the brightening part
smoothed before it is added so that it cannot create a new high frequency edge,
\begin{equation}
\mathbf{v} = \big(\,\mathrm{smooth}(\mathrm{ReLU}(\boldsymbol{\alpha}) \odot \mathbf{b})
 + (\boldsymbol{\alpha} - \mathrm{ReLU}(\boldsymbol{\alpha})) \odot \mathbf{b}\,\big) \odot \mathbf{g},
\label{eq:apply}
\end{equation}
where $\mathrm{smooth}$ is a three by three mean filter. The darkening part is left unsmoothed because
it has zero mean and cancels itself. In the background, defined by the complement of the gate eroded
by two pixels, a masked three by three mean filter is blended in at half strength, which lowers the
background standard deviation without letting blur cross into tissue. Finally the edge proposal is
added and the result is clipped to the valid range,
\begin{equation}
\mathbf{q} = \mathrm{clamp}\big(\mathbf{b} + \mathbf{v} + \boldsymbol{\delta}_{\mathrm{e}} \odot \mathbf{g},\, 0,\, 1\big).
\label{eq:candidate}
\end{equation}

\subsection{Stage six, the safety decision}
\label{sec:safety}
The candidate $\mathbf{q}$ is not returned directly. Three constraints are checked on it.

\textbf{Aggregate quality must not fall.} Define for any image $\mathbf{t}$ the energy
\begin{equation}
E(\mathbf{t}) = 1 - \tfrac{1}{6}\sum_{i=1}^{6} s_i(\mathbf{t})
 + 0.1 \sum_{i=1}^{6} \max\big(0,\, 0.5 - s_i(\mathbf{t})\big),
\label{eq:energy}
\end{equation}
which is small when the image is good, and whose second term penalises any single property below one
half, so a good average cannot hide one broken property. The first constraint requires
$E(\mathbf{q}) \le E(\mathbf{b}) + 0.10$.

\textbf{No property may be sacrificed for another.} The second constraint requires that at least one
property improves by more than $0.01$ and that no property falls by more than $0.10$.

\textbf{The image must not move far.} The third constraint requires
$\lVert \mathbf{q} - \mathbf{b} \rVert_2 \,/\, \lVert \mathbf{b} \rVert_2 < 0.2$.

Let $n \in \{0,1,2,3\}$ be the number of constraints that hold. The output is
\begin{equation}
\hat{\mathbf{x}} = (1 - w)\,\mathbf{b} + w\,\mathbf{q},
\qquad
w = \begin{cases}
1, & n \ge 2, \\[2pt]
0.05 + 0.10\,\dfrac{n}{3}, & n < 2 .
\end{cases}
\label{eq:blend}
\end{equation}
Otherwise the output falls back to the backbone while retaining five to eight percent of the proposal,
so the decision remains visible rather than silently discarded.


\begin{figure*}[!tb]
\centering
\includegraphics[width=0.783\textwidth]{figures/architecture.pdf}
\caption{The six stages. A frozen backbone and a small head produce the image and the cooperation
map. Two networks propose changes. Six properties are measured. Five rules allocate. The change is
bounded and gated. Three constraints decide how much is kept.}
\label{fig:architecture}
\end{figure*}


\begin{table*}[!tb]
\centering
\caption{The six clinical properties in closed form. Here $e$ is the Sobel edge magnitude of
$\mathbf{b}$, $\bar e$ its mean, $b_e = \mathbf{1}[e > \bar e]$, $\sigma_k$ and $\operatorname{Var}_k$
the local standard deviation and variance over a $k \times k$ window, $\bar\sigma$ the mean of
$\sigma_9$, $v(\cdot)$ the depth profile averaged across width, $\nabla_z$ the derivative along depth,
$\rho_h$ the one pixel horizontal autocorrelation, $\operatorname{sigm}$ the logistic function,
$\varepsilon$ a small constant, and $cv$ the local coefficient of variation of the log residual.}
\label{tab:predicates}
\setlength{\tabcolsep}{3pt}
\scriptsize
\begin{tabular}{clp{0.26\textwidth}p{0.225\textwidth}p{0.245\textwidth}c}
\toprule
$P_i$ & Property & Score $s_i$ & Map $\mathbf{m}_i$ & Where the map is large & $\tau_i$ \\
\midrule
$P_1$ & edge continuity &
$0.4\,c_{\text{cont}}+0.3\,c_{\text{efr}}+0.3\,c_{\text{ang}}$ &
$(1-\tilde e)(1-b_e)$, \ $\tilde e=e/(\max e+\varepsilon)$ &
off detected edges, where a boundary should continue but does not & 0.60 \\
$P_2$ & local contrast &
$0.20\,s_{\text{range}}+0.35\,s_{\text{cnr}}+0.20\,s_{\text{var}}+0.25\,s_{\text{cov}}$ &
$\big[\operatorname{ReLU}(\bar\sigma-\sigma_9(\mathbf{b}))/(\bar\sigma+\varepsilon)+0.3\,\mathbf{1}[\sigma_9(\mathbf{b})<0.01]\big]_{[0,1]}$ &
where $\sigma_9$ falls below the image mean, and in nearly flat tissue & 0.50 \\
$P_3$ & flat smoothness &
$0.6\,q_{\text{var}}+0.4\,q_{\text{res}}$ &
$F \odot \sigma_5(\mathbf{b})/(\max \sigma_5(\mathbf{b})+\varepsilon)$ &
inside flat regions that still carry residual variation & 0.75 \\
$P_4$ & structural coherence &
$0.4\,s_{\text{self}}+0.3\,s_{\text{reg}}+0.3\,s_{\text{vert}}$ &
$e(e)/(\max e(e)+\varepsilon)$ &
where the edge map is itself rough, a sign of broken structure & 0.50 \\
$P_5$ & speckle conformity &
$0.7\,\operatorname{mean}(\exp(-|cv-cv^\star|/0.15))+0.3\,\exp(-5|\rho_h(z)|)$ &
$1-\exp(-|cv-cv^\star|/0.15)$ &
where $cv$ departs from the value expected for OCT speckle & 0.30 \\
$P_6$ & layer visibility &
$0.4\,s_{\text{clr}}+0.3\,s_{\text{dep}}+0.3\,s_{\text{smo}}$ &
$\operatorname{Var}_7(e)/(\max \operatorname{Var}_7(e)+\varepsilon)$ &
where the local variance of the edge map is high across depth & 0.70 \\
\midrule
\multicolumn{6}{p{0.97\textwidth}}{
\textbf{Subterms.}
$P_1$.
$c_{\text{cont}}=\langle\operatorname{erode}(\operatorname{dilate}(b_e)),b_e\rangle/(\|b_e\|_1+\varepsilon)$,
$c_{\text{efr}}=\exp(-5|r_e-0.2|)$ with
$r_e=\|\mathbf{1}[e>1.5\bar e]\|_1/(\|\mathbf{1}[e<0.5\bar e]\|_1+\varepsilon)$,
and $c_{\text{ang}}=\exp(-2\,\operatorname{mean}(\sigma_5(\theta)))$ with $\theta$ the Sobel edge angle of $\mathbf{b}$.
$P_2$.
$s_{\text{range}}=\operatorname{sigm}(10(\Delta_r-0.3))$ with $\Delta_r=\max(\mathbf{b})-\min(\mathbf{b})$,
$s_{\text{cnr}}=\operatorname{sigm}(15(\bar\sigma-0.05))$,
$s_{\text{var}}=\max(0.2,\operatorname{sigm}(-20(v_\sigma-0.02)+2))$ with $v_\sigma=\operatorname{Var}(\sigma_9(\mathbf{b}))$,
and $s_{\text{cov}}=\operatorname{sigm}(5(\gamma-0.4))$ with $\gamma=\operatorname{mean}(\mathbf{1}[\sigma_9(\mathbf{b})>0.3\bar\sigma])$.
$P_3$.
$F=\mathbf{1}[e<0.3\bar e]$,
$q_{\text{var}}=1-\langle F,\sigma_5(\mathbf{b})/(\sigma_5(\mathbf{y})+\varepsilon)\rangle/(\|F\|_1+\varepsilon)$,
and $q_{\text{res}}=\exp(-5\,\operatorname{mean}|\mathbf{y}-\mathbf{b}|)$.
$P_4$.
$s_{\text{self}}$ is the mean Pearson correlation between $\mathbf{b}$ and its shifts by $(0,2)$, $(2,0)$ and $(2,2)$ pixels,
$s_{\text{reg}}=\exp(-12\,\operatorname{mean}(e(e)))$,
and $s_{\text{vert}}=\operatorname{clip}(\max |\nabla_z v(\mathbf{b})|/(8\,\operatorname{mean}|\nabla_z v(\mathbf{b})|),0,1)$.
$P_5$.
$z=\log(\mathbf{y}+\varepsilon)-\log(\mathbf{b}+\varepsilon)$,
$cv=\sigma_{15}(|z|)/(\mu_{15}(|z|)+\varepsilon)$ with $\mu_{15}$ the local mean,
and $cv^\star(d)$ is the target coefficient of variation of OCT speckle at axial depth $d$ under the Gamma K model of the released code.
$P_6$.
$s_{\text{clr}}=\operatorname{mean}(\operatorname{clip}(\max |\nabla_z v(\mathbf{b})|/(10\,\operatorname{mean}|\nabla_z v(\mathbf{b})|),0,1))$,
$s_{\text{dep}}=\operatorname{clip}(20\,\operatorname{Var}(v(\mathbf{b})),0,1)$,
and $s_{\text{smo}}=1-\operatorname{clip}(10\,\operatorname{mean}(\operatorname{Var}_z(\operatorname{mean}_w e)),0,1)$
with $\operatorname{mean}_w$ the mean across width and $\operatorname{Var}_z$ the variance along depth.} \\
\bottomrule
\end{tabular}
\end{table*}
